% !TeX root = ../深度学习与大语言模型技术原理.tex
% 由主文档通过 \input 载入；请编译主文档。
\setcounter{chapter}{14} % 保留原书章号。
\chapter{检索增强生成与复杂检索}
\label{ch:rag}

\section{从参数知识到可追溯证据}

\emph{检索增强生成（retrieval-augmented generation, RAG）}把外部文档作为条件，用来生成有依据的回答。一个概率化视角为 $p(y\mid q)=\sum_d p_\eta(d\mid q)p_\theta(y\mid q,d)$；实际系统通常只取若干候选，并把它们拼入上下文，而非精确边缘化整个文档库\cite{rag}。检索结果决定信息来源，生成模型负责组织与推理，两部分的错误需要分开定位。

文档切块应保留段落、标题、表格说明和来源位置。块过短会失去条件，块过长会降低匹配精度并消耗上下文。重叠块减少边界断裂，却增加重复与检索冗余。索引单元可以较小，传给模型的证据单元可以扩展到父段落；二者不必相同。

引用需要把最终陈述关联到真实可访问的证据位置。检索到文档并附一个链接，并不证明该文档支持附近结论。应区分引用可解析、引用内容相关和引用确实蕴含陈述三个层次；权限过滤应在检索和返回阶段生效，避免先召回无权文档再寄希望于模型不泄露。

\section{查询改写、分解与 HyDE}

查询改写（query rewriting）把省略、口语或多轮指代表达转成更适合检索的形式；查询扩展（query expansion）增加同义词、别名和相关术语。多查询检索用若干表达分别召回后合并。改写若改变否定、时间、主体或数字条件，会把搜索引向错误问题，因此应保留原始约束并比较召回增益。

查询分解（query decomposition）把复杂问题拆为子问题。若后一个问题需要前一结果，则不能机械并行。例如比较两家公司某年的指标，可以分别检索两者再对齐口径；查找某奖项得主的导师，则需先确定得主。子查询之间的依赖构成有向图，执行与停止应围绕依赖关系管理。

假设文档嵌入（Hypothetical Document Embeddings, HyDE）先生成一个可能包含答案的假设文档，再编码它检索真实文档\cite{hyde}。其动机是让查询更接近文档表达空间；生成的假设文本只是检索媒介，可能包含虚构事实，不能充当答案证据。最终回答仍需回到实际召回来源。

\section{上下文化索引与迭代检索}

上下文化检索（contextual retrieval）在索引前为切块补充父文档语境，例如公司、年份、章节和实体。孤立句“增长20\%”缺少基准，附加真实元数据后更容易被正确检索。若补充语境由模型生成，必须可回溯到原文，避免在索引阶段把幻觉固化成长期知识\cite{contextual}。

智能体式检索增强生成（agentic RAG）根据当前证据决定是否继续检索：识别缺失事实，产生新查询，检查结果，直到证据充分或预算用尽。它比一次检索多了一个决策层，也多了查询漂移、重复搜索和错误自信的风险。证据充分性应按任务需要判断，例如比较题必须覆盖两个对象和相同时间口径，而不是简单以文档数量为准。

停止规则应同时包含成功条件和资源上限。达到预算但证据不足时，应保留未解决项；不能把“没有搜到反例”视为结论已证实。评估迭代检索时，需要记录每次检索的增量证据和调用成本，检查多轮是否确实补足了缺失信息。

\section{图检索与全局总结}

图检索增强生成（Graph RAG）把实体及关系组织为图，可沿邻接关系、路径或社区摘要检索。多跳问题能通过路径定位相关证据，跨文档总结能利用社区级组织；微软提出的 GraphRAG 研究包含面向全局问题的社区摘要路线\cite{graphrag}，并非所有冠以 Graph RAG 的实现都相同。

图中的边可能来自人工结构化数据，也可能由模型抽取。后一种情况下，实体消歧、关系方向、时间和来源必须保留，抽取边不能自动被当作事实。图更新还涉及删除过时关系、版本一致性和摘要重建，比普通向量插入更复杂。

如果任务只需在短文档中查一个精确句子，图构建未必值得；如果需要跨文档关系和全局组织，图可能有帮助。应以同一数据、同一预算比较普通混合检索、迭代检索和图方案，并分析收益来自结构还是额外索引时计算。

\section{检索与生成的分层指标}

设查询的已标注相关文档集合为 $R$，前 $K$ 个结果为 $S_K$，则 $\operatorname{Recall@K}=|R\cap S_K|/|R|$、$\operatorname{Precision@K}=|R\cap S_K|/K$，假定 $|R|>0$ 且返回 $K$ 项。相关性标注不完整时，这些指标会受到未标注文档的影响。平均倒数排名（MRR）平均第一个相关结果排名的倒数，无相关结果取零。

折损累计增益 $\operatorname{DCG@K}=\sum_{i=1}^K(2^{\mathrm{rel}_i}-1)/\log_2(i+1)$ 允许分级相关性，归一化为 $\operatorname{NDCG@K}=\operatorname{DCG@K}/\operatorname{IDCG@K}$，理想增益为零时需规定评测约定。不同指标偏重覆盖、头部命中或排序质量，不能互相替代。

生成侧测量正确性、完整性、忠实性与引用支持率。忠实于错误文档仍可能事实错误，事实碰巧正确也可能没有证据支持。定位问题时可分别给模型真实检索证据和人工理想证据：若理想证据下仍失败，应优先检查生成与任务理解；若理想证据显著改善，再分析召回、重排和上下文组织。

